\documentclass[10pt,journal,compsoc,a4paper]{IEEEtran}

\usepackage[utf8]{inputenc}
\usepackage[T1]{fontenc}
\usepackage{amsmath,amsfonts,amssymb}
\usepackage{graphicx}
\usepackage{booktabs}
\usepackage{multirow}
\usepackage{array}
\usepackage{tabularx}
\usepackage{xcolor}
\usepackage{hyperref}
\usepackage{microtype}
\usepackage{cite}
\usepackage{siunitx}
\usepackage{float}

\hypersetup{
    colorlinks=true,
    linkcolor=blue!70!black,
    citecolor=red!70!black,
    urlcolor=blue!80!black
}

\begin{document}

\title{Relatório Técnico de Validação Experimental e Avaliação Multidimensional: Arquitetura xApp RDL (Fase 1: H-RDL) vs. Baseline no Near-RT RIC O-RAN}

\author{
    \IEEEauthorblockN{George Barbosa\IEEEauthorrefmark{1}}\\
    \IEEEauthorblockA{\IEEEauthorrefmark{1}Pesquisa e Engenharia Open RAN / Near-RT RIC\\
    Repositório GitHub: \href{https://github.com/georgebarbosa3090/XApp-RDL-F1}{https://github.com/georgebarbosa3090/XApp-RDL-F1}\\
    Notebook Interativo Google Colab: \href{https://colab.research.google.com/github/georgebarbosa3090/XApp-RDL-F1/blob/main/notebooks/rdl_colab_scikit_learn.ipynb}{rdl\_colab\_scikit\_learn.ipynb}\\
    Data da Execução: 31 de Agosto de 2026}
}

\markboth{Relatório Técnico de Engenharia e Pesquisa O-RAN --- 31 de Agosto de 2026}{Barbosa: Validação Experimental xApp RDL Fase 1}

\IEEEtitleabstractindextext{
\begin{abstract}
Este relatório técnico apresenta os resultados quantitativos, empíricos e reproduzíveis obtidos na campanha experimental executada em \textbf{31 de Agosto de 2026} referente à arquitetura \textbf{xApp RDL (Resource and Decision Layer) --- Fase 1 (H-RDL Determinística)}. A avaliação confronta o regime operacional desregulador (\textit{Baseline Sem RDL}), no qual múltiplas xApps concorrentes (\textit{xSlice}, \textit{Energy Saving} e \textit{Traffic Steering}) disputam os recursos de rádio na Near-RT RIC sem mediação, contra o regime de governança orquestrada pela H-RDL conectada via interfaces O-RAN E2 (E2SM-KPM / E2SM-RC). A co-simulação foi conduzida no simulador de redes 5G-NR ns-3 v3.40 (módulos 5G-LENA e NORI) na faixa de 3.5 GHz (Banda n78) acoplado a um cluster Kubernetes Near-RT RIC. Adicionalmente, um pipeline de Machine Learning (Scikit-Learn) com validação cruzada estratificada em 10 folds e modelos ensemble foi avaliado para classificação preditiva de conflitos inter-xApps. Os dados evidenciam que a introdução da RDL reduziu a latência média URLLC em \textbf{75.8\%} (de 11.79 ms para 2.85 ms), reduziu o percentil 99 da latência em \textbf{97.8\%} (de 139.41 ms para 3.09 ms), erradicou totalmente as violações de SLA URLLC (\textbf{0.0\%} de estouro), elevou a Taxa de Entrega de Pacotes (PDR) de 39.28\% para \textbf{99.53\%}, mitigou \textbf{97.8\%} dos conflitos inter-xApp em tempo de decisão médio de \textbf{14.2 ms} (meta $<$ 50 ms) e propiciou ganho de eficiência energética de \textbf{+14.5\%} (Bits/Joule).
\end{abstract}

\begin{IEEEkeywords}
Open RAN, Near-RT RIC, xApp, RDL, Resource Allocation, QoS/SLA, Machine Learning, ns-3, 5G-LENA, Conflitos de Controle.
\end{IEEEkeywords}
}

\maketitle
\IEEEdisplaynontitleabstractindextext
\IEEEpeerreviewmaketitle

\section{Introdução e Contexto O-RAN}
\IEEEPARstart{A}{arquitetura} Open RAN (O-RAN Alliance) descentraliza a pilha de protocolo de acesso via rádio através de interfaces abertas e controlabilidade programável orientada por microserviços inteligentes (\textit{xApps}) executados no \textit{Near-Real-Time RAN Intelligent Controller} (Near-RT RIC)~\cite{oran_wg1, oran_wg3}. Embora a modularidade de xApps independentes permita a especialização funcional de fatiamento de rede (\textit{Network Slicing}), economia energética (\textit{Energy Saving}) e direcionamento de tráfego (\textit{Traffic Steering}), a operação concorrente e não coordenada dessas aplicações induz a graves conflitos de controle (\textit{Cross-xApp Control Conflicts})~\cite{bonati2021openran, polese2023understanding}.

Tais conflitos ocorrem em três eixos principais:
\begin{enumerate}
    \item \textbf{Conflito de Recursos de Frequência (PRB):} A xApp de fatiamento reserva blocos de recursos físicos (\textit{Physical Resource Blocks} - PRBs) prioritários para Ultra-Reliable Low-Latency Communication (URLLC), enquanto picos de tráfego de Enhanced Mobile Broadband (eMBB) sobrecarregam o escalonador MAC;
    \item \textbf{Conflito de Potência e Cobertura:} A xApp de economia de energia atenua a potência de transmissão ($P_{\text{TX}}$) dos gNodeBs, degradando a relação sinal-ruído mais interferência (SINR) de fluxos críticos de missão;
    \item \textbf{Conflito de Mobilidade e Handover:} A xApp de balanceamento de tráfego força handovers imediatos que colidem com os SLAs estritos de latência, gerando oscilações espúrias (\textit{Handover Ping-Pong}).
\end{enumerate}

Para solucionar essa lacuna fundamental de governança, o projeto \textbf{xApp RDL (Resource and Decision Layer)} foi concebido como uma camada intermediária de arbitragem e conformidade de políticas em tempo real (Near-RT: $10\text{ ms} - 1\text{ s}$). A \textbf{Fase 1 (H-RDL)} estabelece uma abordagem heurística determinística com barreiras de segurança (\textit{Safety Guards}) e priorização hierárquica baseada em 3GPP 5G QoS Identifiers (5QI).

\section{Metodologia Experimental e Configuração do Ambiente}
A campanha experimental executada em 31 de Agosto de 2026 empregou uma bancada de co-simulação integrada de alta fidelidade:
\begin{itemize}
    \item \textbf{Simulador de Rede 5G NR:} ns-3 v3.40 integrado com os frameworks 5G-LENA (CTTC) e NORI (OpenRAN ns-3 E2 Interface);
    \item \textbf{Topologia de Rádio:} 1 gNodeB gNB central e 34 UEs heterogêneos distribuídos geograficamente, cobrindo três classes de tráfego concorrente: URLLC (5QI 82), eMBB (5QI 9) e mMTC (5QI 79);
    \item \textbf{Parâmetros de RF:} Banda n78 (frequência central de 3.5 GHz, largura de banda de 50 MHz, numerologia $\mu=1$, espaçamento de subportadoras $\Delta f = 30\text{ kHz}$);
    \item \textbf{Controlador Near-RT RIC:} Cluster Kubernetes k3d com namespace dedicado \texttt{ricxapp}, integrado via terminação E2 SCTP (porta 36422), barramento RMR (portas 4560/4561) e métricas Prometheus (porta 8081).
\end{itemize}

O protocolo de teste foi dividido em dois cenários estritamente controlados sob o mesmo perfil de carga e mobilidade:
\begin{enumerate}
    \item \textbf{Cenário Baseline (Sem RDL):} Três xApps de referência (\textit{xSlice}, \textit{Energy Saving}, \textit{Traffic Steering}) enviam comandos simultâneos diretamente ao gNodeB via E2SM-RC sem coordenação;
    \item \textbf{Cenário Fase 1 (H-RDL):} Todas as propostas de ação passam pela xApp RDL, que aplica verificação de conflitos em matriz de impacto cruzado, impondo limites rígidos de SLA antes de enviar comandos de controle via E2.
\end{enumerate}

\section{Resultados Experimentais de Rede e Governança O-RAN}
A Tabela~\ref{tab:metricas_rede} resume o comparativo multidimensional das métricas de desempenho coletadas e consolidadas em 31 de Agosto de 2026.

\begin{table*}[t]
\centering
\caption{Avaliação Comparativa Multidimensional: Baseline (Sem RDL) vs. Fase 1 (H-RDL) --- Execução de 31 de Agosto de 2026}
\label{tab:metricas_rede}
\begin{tabularx}{\textwidth}{l l c c c X}
\toprule
\textbf{Domínio de Avaliação} & \textbf{Métrica Científica} & \textbf{Baseline (Sem RDL)} & \textbf{Fase 1: H-RDL} & \textbf{Variação / Ganho} & \textbf{Impacto Técnico no 5G/O-RAN} \\
\midrule
\multirow{4}{*}{\textbf{QoS \& Latência URLLC}} 
& Latência Média URLLC & \SI{11.79}{\milli\second} & \textbf{\SI{2.85}{\milli\second}} & \textbf{-75.8\%} & Redução drástica de filas na camada MAC \\
& Latência Percentil 95 (P95) & \SI{53.14}{\milli\second} & \textbf{\SI{3.08}{\milli\second}} & \textbf{-94.2\%} & Estabilidade de cauda determinística \\
& Latência Percentil 99 (P99) & \SI{139.41}{\milli\second} & \textbf{\SI{3.09}{\milli\second}} & \textbf{-97.8\%} & Garantia estrita de requisitos 3GPP \\
& Violação de SLA URLLC ($>5\text{ms}$) & 29.17\% & \textbf{0.0\%} & \textbf{-100.0\%} & Eliminação completa de estouro de SLA \\
\midrule
\multirow{2}{*}{\textbf{Confiabilidade \& Perda}} 
& Taxa de Entrega de Pacotes (PDR) & 39.28\% & \textbf{99.53\%} & \textbf{+153.4\%} & Erradicação de perdas por contenção \\
& Taxa de Perda de Pacotes (PLR) & 60.72\% & \textbf{0.47\%} & \textbf{-99.2\%} & Queda expressiva de retransmissões HARQ \\
\midrule
\multirow{2}{*}{\textbf{Throughput \& Equidade}} 
& Throughput Médio por Fluxo & \SI{2.17}{\mega\bit\per\second} & \textbf{\SI{37.04}{\mega\bit\per\second}} & \textbf{+1606.9\%} & Vazão útil destravada com escalonamento justo \\
& Índice de Equidade de Jain & 0.1414 & \textbf{0.9164} & \textbf{+548.1\%} & Coexistência harmônica inter-slice \\
\midrule
\multirow{5}{*}{\textbf{Governança \& Conflitos}} 
& Taxa de Ocorrência de Conflitos & 34.67\% & 30.67\% & -11.5\% & Demanda natural de tráfego concorrente \\
& Conflitos Não Mitigados (\%) & 34.67\% & \textbf{0.67\%} & \textbf{-98.1\%} & Quase anulação de colisões de controle \\
& Eficiência de Arbitragem RDL & 0.0\% & \textbf{97.83\%} & \textbf{+97.8 p.p.} & Resolução proativa via Safety Guards \\
& Latência de Decisão da RDL & N/A & \textbf{\SI{14.2}{\milli\second}} & Meta $<50\text{ms}$ & Total conformidade com Near-RT RIC \\
& Handover Ping-Pong & \SI{22.0}{ev\per\minute} & \textbf{\SI{0.0}{ev\per\minute}} & \textbf{-100.0\%} & Estabilidade absoluta de mobilidade \\
\midrule
\multirow{3}{*}{\textbf{Eficiência Energética}} 
& Índice Bits/Joule Normalizado & 1.00x & \textbf{1.145x} & \textbf{+14.5\%} & Otimização sustentável de potência TX \\
& Potência Média de Transmissão & \SI{39.01}{\decibel m} & \textbf{\SI{34.87}{\decibel m}} & \textbf{-4.14 dBm} & Minimização de interferência de célula vizinha \\
& Conformidade Global de SLA & 65.33\% & \textbf{100.0\%} & \textbf{+34.7 p.p.} & Satisfação integral das fatias de rede \\
\bottomrule
\end{tabularx}
\end{table*}

\subsection{Análise de Latência e Confiabilidade URLLC}
Como demonstrado na Tabela~\ref{tab:metricas_rede}, o regime Baseline apresentou grave degradação de cauda, atingindo latência P99 de \SI{139.41}{\milli\second}, em virtude do esgotamento de PRBs causado pelas requisições desordenadas da xApp eMBB. Sob o controle da H-RDL, a latência média foi comprimida para \SI{2.85}{\milli\second}, e o P99 permaneceu contido em \SI{3.09}{\milli\second}, sem nenhuma violação do limiar de \SI{5}{\milli\second} estipulado pelo 3GPP TS 22.261.

\subsection{Eficiência de Governança e Arbitragem RDL}
A taxa de conflitos potenciais no sistema permaneceu em torno de 30--35\%, o que comprova que a carga e a disputa de recursos foram idênticas em ambos os cenários. No entanto, enquanto o Baseline deixou \textbf{34.67\%} de conflitos ativos (resultando em oscilações de 22 handovers/minuto), a H-RDL atingiu eficiência de arbitragem de \textbf{97.83\%}, resolvendo as disputas em tempo médio de \textbf{\SI{14.2}{\milli\second}}, margem de folga substancial em relação ao limite estipulado pelo Near-RT RIC ($<50\text{ ms}$).

\section{Benchmark de Machine Learning para Classificação de Conflitos}
Com o propósito de viabilizar a transição para a Fase 2 (arbitragem proativa baseada em inteligência artificial), foi desenvolvido um pipeline de Machine Learning calibrado com os atributos de telemetria de rádio gerados nas simulações. 

\subsection{Engenharia de Atributos (\textit{Feature Engineering})}
Foram calculados os seguintes atributos matemáticos a partir das variáveis primitivas:
\begin{align}
\text{SINR}_{\text{linear}} &= 10^{\frac{\text{SINR}_{\text{dB}}}{10}} \\
\text{CapacityProxy} &= \log_2(1 + \text{SINR}_{\text{linear}}) \\
\text{PowerPerPRB} &= \frac{10^{\frac{P_{\text{TX,dBm}}}{10}}}{\text{PRB}_{\text{demanded}} + \epsilon} \\
\text{StressIndex} &= \left(\frac{\text{Load}_{\text{Mbps}}}{100}\right) \times \left(\frac{\text{PRB}_{\text{demanded}}}{272}\right)
\end{align}

\subsection{Avaliação e Comparação de Modelos de ML}
Seis modelos foram submetidos a validação cruzada estratificada em 10 folds (\textit{Stratified 10-Fold CV}) e avaliados em conjunto de teste independente (75 amostras). Os resultados são exibidos na Tabela~\ref{tab:ml_benchmark}.

\begin{table*}[t]
\centering
\caption{Benchmark Científico de Algoritmos de Machine Learning para Detecção de Conflitos O-RAN (Notebook Google Colab)}
\label{tab:ml_benchmark}
\begin{tabularx}{\textwidth}{l c c c c c c c c c}
\toprule
\textbf{Algoritmo} & \textbf{CV Accuracy} & \textbf{CV F1-Score} & \textbf{CV ROC-AUC} & \textbf{Test Acc (\%)} & \textbf{Test Prec (\%)} & \textbf{Test Rec (\%)} & \textbf{Test F1} & \textbf{Test ROC-AUC} & \textbf{Brier Score} \\
\midrule
Decision Tree & 95.99\% $\pm$ 3.12\% & 0.9362 $\pm$ 0.0519 & 0.9548 $\pm$ 0.0423 & 97.33 & 95.83 & 95.83 & 0.9583 & 0.9788 & 0.0219 \\
Random Forest (Tuned) & 97.33\% $\pm$ 3.56\% & 0.9547 $\pm$ 0.0617 & 0.9972 $\pm$ 0.0060 & \textbf{100.00} & \textbf{100.00} & \textbf{100.00} & \textbf{1.0000} & \textbf{1.0000} & 0.0064 \\
Extra Trees & 94.68\% $\pm$ 5.16\% & 0.9148 $\pm$ 0.0843 & 0.9901 $\pm$ 0.0165 & \textbf{100.00} & \textbf{100.00} & \textbf{100.00} & \textbf{1.0000} & \textbf{1.0000} & 0.0341 \\
Gradient Boosting & 95.53\% $\pm$ 4.02\% & 0.9267 $\pm$ 0.0707 & 0.9972 $\pm$ 0.0060 & 98.67 & \textbf{100.00} & 95.83 & 0.9787 & 0.9992 & 0.0129 \\
HistGradientBoosting & \textbf{98.66\% $\pm$ 2.05\%} & \textbf{0.9790 $\pm$ 0.0322} & \textbf{1.0000 $\pm$ 0.0000} & \textbf{100.00} & \textbf{100.00} & \textbf{100.00} & \textbf{1.0000} & \textbf{1.0000} & \textbf{0.0000} \\
\textbf{Ensemble (RF+ET+GB+HGB)} & 97.31\% $\pm$ 3.57\% & 0.9547 $\pm$ 0.0617 & 0.9990 $\pm$ 0.0029 & \textbf{100.00} & \textbf{100.00} & \textbf{100.00} & \textbf{1.0000} & \textbf{1.0000} & 0.0051 \\
\bottomrule
\end{tabularx}
\end{table*}

Os resultados do modelo ensemble e do HistGradientBoosting demonstram que as características de estresse de tráfego (\texttt{stress\_index}) e densidade de potência por PRB (\texttt{power\_per\_prb}) fornecem separabilidade quase perfeita para prever colisões de decisão antes da injeção de comandos na RAN.

\section{Reprodutibilidade e Artefatos de Software}
Todos os experimentos descritos neste relatório foram versionados e disponibilizados publicamente para auditoria e reprodutibilidade:
\begin{itemize}
    \item \textbf{Código-fonte e Helm Charts:} \href{https://github.com/georgebarbosa3090/XApp-RDL-F1}{https://github.com/georgebarbosa3090/XApp-RDL-F1}
    \item \textbf{Caderno Interativo Google Colab:} \href{https://colab.research.google.com/github/georgebarbosa3090/XApp-RDL-F1/blob/main/notebooks/rdl_colab_scikit_learn.ipynb}{rdl\_colab\_scikit\_learn.ipynb}
    \item \textbf{Diretório de Resultados Diários:} \texttt{experiments/results/2026-08-31/run\_103501/}
    \item \textbf{Datasets de Telemetria:} \texttt{dataset\_flow\_metrics.csv} e \texttt{dataset\_rdl\_decisions\_ml.csv}.
\end{itemize}

\section{Conclusão e Próximos Passos}
A validação experimental de 31 de Agosto de 2026 comprovou conclusivamente que a introdução da xApp RDL estabelece uma governança robusta e determinística sobre o Near-RT RIC. A solução garantiu \textbf{100.0\%} de conformidade com os SLAs de fatiamento de rede, eliminou eventos destrutivos de handover ping-pong e reduziu a latência URLLC para \textbf{\SI{2.85}{\milli\second}}, com overhead de decisão de apenas \textbf{\SI{14.2}{\milli\second}}. 

Como próximo passo evolutivo, a \textbf{Fase 2 (CA-RDL / MARL)} incorporará o modelo preditivo aqui validado a um algoritmo de Multi-Agent Reinforcement Learning (MAPPO/Actor-Critic), dotando a RDL de capacidade de aprendizado contextual autônomo.

\begin{thebibliography}{10}
\bibitem{oran_wg1}
O-RAN Alliance WG1, ``O-RAN Architecture Description 11.00,'' \emph{O-RAN Technical Specification}, Jul. 2024.

\bibitem{oran_wg3}
O-RAN Alliance WG3, ``Near-Real-Time RAN Intelligent Controller Architecture \& E2 Service Models (E2SM-KPM / E2SM-RC),'' \emph{O-RAN Technical Specification}, 2023.

\bibitem{bonati2021openran}
L.~Bonati, M.~Polese, S.~D'Oro, S.~Basagni, and T.~Melodia, ``OpenRAN Gym: An Open, Interoperable, and Experimental Framework for O-RAN Research,'' \emph{IEEE Transactions on Wireless Communications}, vol.~22, no.~6, pp.~3981--3995, 2023.

\bibitem{polese2023understanding}
M.~Polese, L.~Bonati, S.~D'Oro, and T.~Melodia, ``Understanding O-RAN: Architecture, Interfaces, Algorithms, Security, and Research Challenges,'' \emph{IEEE Communications Surveys \& Tutorials}, vol.~25, no.~2, pp.~1376--1411, 2023.

\bibitem{3gpp_ts22261}
3GPP, ``Service requirements for the 5G system; Stage 1,'' \emph{3GPP TS 22.261 version 18.5.0 Release 18}, Apr. 2024.
\end{thebibliography}

\end{document}
